% ===== 8 工程实现与现场部署 =====
\section{工程实现与现场部署}
\label{sec:deployment}

前两章已经分别回答“如何验证方法”和“当前证据能支持哪些结论”。对于竞赛现场，另一个同样关键的问题是：在有限时间、有限操作次数和不确定环境下，如何把相机、尺度、实例分割、筛分映射和结果导出组织成一个稳定的操作流程。本文的工程目标不是追求单个模块的极限吞吐，而是保证从设备就位到最终报告之间的每一步都有明确输入、可检查输出和失败后的恢复路径，使现场 30~min 约束能够被端到端流程而不是单一 GPU 推理时间验证。

\subsection{部署目标与系统边界}
\label{sec:deployment-goals}

现场系统需要同时满足四个工程目标。第一是\textbf{几何可追溯}：所有毫米级结果必须能够追溯到本次拍摄的尺度标定，而不能沿用未知来源的固定 mm/px。第二是\textbf{视觉可质检}：实例掩膜、轴长叠加和异常高亮必须能够快速目检，使严重粘连、漏分割或尺度错误在提交前暴露。第三是\textbf{参数可冻结}：比赛正式测试时不应边看结果边重新训练或反复调节下游参数；参数应来自赛前开发/校准阶段并在测试前固定。第四是\textbf{结果可复现}：场景图像、尺度、模型名、$\boldsymbol{\theta}$、$\gamma$、逐颗粒表和场景汇总应一起保存，避免最终只留下无法追溯的截图。

基于这些目标，本文将现场流程划分为三层。\emph{采集层}负责照明、视角、取景和尺度标记；\emph{计算层}负责实例分割、几何测量和筛分/形貌/异常分析；\emph{质检与输出层}负责检查中间结果并生成提交口径。三层之间保持清晰的数据接口：采集层输出图像和尺度，计算层输出实例掩膜与结构化表格，质检层只基于已生成结果判断是否需要重拍或重新运行固定配置，而不在最后一步临时改变方法定义。

\subsection{硬件配置与成像原则}
\label{sec:deployment-hardware}

表\ref{tab:deployment-hardware}给出适合本方案的推荐硬件配置。表中的参数是\textbf{工程建议而非唯一型号要求}：真正影响算法的不是某一品牌，而是图像覆盖范围、空间分辨率、运动模糊、畸变、照明均匀性以及 GPU 是否能在时间预算内完成推理。因此，若现场使用其他设备，应优先保证这些可测量条件，而不是机械复制型号。

\begin{table}[htbp]
\centering
\caption{现场部署的推荐硬件及其与算法误差的关系。}
\label{tab:deployment-hardware}
\small
\renewcommand{\arraystretch}{1.15}
\begin{tabular}{>{\raggedright\arraybackslash}p{2.1cm}>{\raggedright\arraybackslash}p{4.0cm}>{\raggedright\arraybackslash}p{7.5cm}}
\toprule
模块 & 推荐条件 & 作用与理由 \\
\midrule
相机 & 约 12~MP 或更高；固定曝光；若输送带场景优先全局快门 & 保证 5--40~mm 主粒级具有足够像素覆盖；降低运动造成的轮廓拖影 \\
镜头 & 固定焦距、低畸变；安装后锁紧焦距与高度 & 保证标定后 mm/px 在一次测试内稳定，减少视场变化造成的尺度漂移 \\
支架 & 相机光轴尽量垂直于测量平面，刚性固定 & 降低透视缩放差异；若无法近似正俯视，应增加平面单应校正 \\
背景/承载面 & 深色、哑光、几何平整 & 增大骨料与背景的视觉差异，降低高光与复杂纹理对实例边界的干扰 \\
照明 & 大面积漫射、亮度稳定，避免单侧硬阴影 & 降低接触边缘被阴影连接或高光截断的概率 \\
尺度参考 & 标尺或已知尺寸平面标记；尽量与骨料处于同一测量平面 & 为像素到毫米的转换提供本次场景物理锚点 \\
计算设备 & 支持 CUDA 的 NVIDIA GPU；显存满足 Cellpose-SAM 推理 & 缩短实例分割阶段的 wall-clock time；CPU 可运行但现场余量更小 \\
\bottomrule
\end{tabular}
\end{table}

对于自然堆积赛道，成像准备的核心不是把所有颗粒完全摊成互不接触的实验室单层，而是在不改变比赛场景性质的前提下减少不可恢复遮挡。应优先保证相机正俯、焦平面覆盖目标区域、图像不过曝且没有大面积深阴影。若允许整理样品，推荐使可见表层不过度二层堆叠；若比赛要求保持原堆积状态，则应记录遮挡程度，并把严重遮挡导致的不可见颗粒视为单目视觉方法的边界，而不是通过后处理“补出”未观测颗粒。

尺度参考物必须与被测表面尽量共面。如果标尺明显高于或低于骨料测量平面，即使像素测量本身准确，也会因透视比例不同形成系统尺寸误差。当前实现以单一 mm/px 作为场景尺度，因此更适合相机距离相对大、视场近似平面且光轴接近垂直的布置；若现场透视变形明显，下一阶段应使用多点标记和单应/畸变校正把尺度恢复升级为空间位置相关模型。

\subsection{软件流水线与数据接口}
\label{sec:deployment-software}

现场软件链由两个相对独立的命令行阶段组成。第一阶段由 \texttt{imagegrains} 完成 Cellpose-SAM 实例分割和几何测量，核心输出包括预测实例掩膜、用于目检的 composite 图、逐颗粒像素几何 CSV 以及写入毫米轴长的 \texttt{*\_re\_scaled.csv}。第二阶段由 \texttt{aggregate\_screening} 读取逐颗粒表，执行筛分等效粒径、质量代理、形貌分类、异常分类和汇总报告。把两阶段分开有一个工程优势：若上游实例结果已经确认无误，修改报告格式或重算下游统计时无需重复执行昂贵的分割推理。

当前仓库的典型执行方式如下：
\begin{lstlisting}[language=bash,caption={现场处理的两阶段命令示例。正式比赛应将 resolution 和冻结参数替换为本次标定/赛前校准值。}]
python -m imagegrains \
  --img_dir <scene_dir> \
  --model_dir models \
  --out_dir <measurement_out> \
  --resolution <mm_per_px> \
  --gpu True

python -m aggregate_screening \
  --grains <measurement_out>/*_re_scaled.csv \
  --out_dir <report_out>
\end{lstlisting}

正式部署时，推荐再封装一层场景级脚本或 GUI，把图像路径、尺度、模型版本、$\boldsymbol{\theta}$、$\gamma$ 和输出目录写入同一个配置记录。这样可以避免手工命令中最常见的三类错误：选错上一批次 CSV、误用旧 resolution、或在不同场景之间使用不一致参数。自动化封装的目标不是改变算法，而是降低操作失误概率并缩短现场交互时间。

\subsection{现场尺度标定与采集质检}
\label{sec:field-calibration-qc}

尺度标定应在每次相机高度、焦距或取景发生变化后重新执行。最简单的做法是在测量平面放置已知长度 $L_{\mathrm{ref}}$ 的标尺或标记物，读取其图像长度 $l_{\mathrm{ref}}$，计算
\begin{equation}
 r=\frac{L_{\mathrm{ref}}}{l_{\mathrm{ref}}}\quad \mathrm{mm/px}.
\label{eq:field-resolution}
\end{equation}
该 $r$ 随后作为第\ref{sec:scale-calibration}节的物理尺度输入。若同一图像可读取多个参考长度，应比较不同位置得到的 $r_j$；若离散明显，说明透视或畸变已经不能用单一 mm/px 近似，应先调整相机姿态或启用更完整的几何校正。

图像采集完成后，不应立即进入最终推理，而应先做一次低成本质量检查。建议至少检查五项：目标区域是否完整入镜；参考物是否清晰且共面；是否存在过曝/欠曝；是否有大面积硬阴影；是否因相机振动或输送带运动出现明显模糊。任何一项失败都应优先重拍，因为这些问题一旦进入实例分割，会同时影响边界、轴长和所有下游级配指标，而后处理很难可靠修复。

推理完成后的第二次质检针对模型输出。首先查看 composite 或 detection overview，确认密集接触区域是否发生明显整片 merge，边缘颗粒是否被大量截断，小颗粒是否异常碎裂；其次查看 axes overlay，确认轴长方向与颗粒轮廓一致且比例尺量级合理；最后查看场景报告中的实例数量、异常比例和 $D_p$ 是否存在数量级异常。这里的质检用于识别明显失败，不应演变成“看到 $D_{50}$ 不符合预期就反复调参”，否则会把测试样本信息泄漏回模型选择过程。

\subsection{30~min 现场操作流程与时间预算}
\label{sec:30min-sop}

根据比赛的现场时间约束，本文建议采用表\ref{tab:30min-budget}所示的\textbf{时间预算上限}组织操作。该表不是对当前系统已经完成的正式端到端计时声明，而是将 30~min 拆解为可执行检查点，并为重拍或软件异常预留缓冲。第\ref{sec:runtime-results}节已有约 27~s/张的算法阶段历史记录，因此在 GPU 正常工作的情况下，真正需要工程优化的往往不是单次推理，而是设备布置、尺度确认、人工质检和故障恢复。

\begin{table}[htbp]
\centering
\caption{建议的 30~min 现场 SOP 与预算上限。正式报告应使用现场完整演练的实测时间替换“预算”。}
\label{tab:30min-budget}
\small
\renewcommand{\arraystretch}{1.15}
\begin{tabular}{>{\centering\arraybackslash}p{1.0cm}>{\raggedright\arraybackslash}p{3.1cm}>{\centering\arraybackslash}p{1.8cm}>{\raggedright\arraybackslash}p{7.3cm}}
\toprule
步骤 & 操作 & 预算 & 完成判据 \\
\midrule
1 & 设备就位与视场调整 & 4~min & 相机固定、目标区域完整、照明稳定 \\
2 & 尺度标定 & 3~min & 获得本场景 mm/px；多位置参考无明显冲突 \\
3 & 图像采集与快速质检 & 3~min & 至少一张清晰可用图像，参考物和骨料均无遮挡到无法读取 \\
4 & 实例分割与几何测量 & 5~min & 生成 mask、composite 和 \texttt{re\_scaled.csv}；无大面积 merge \\
5 & 筛分/形貌/异常分析 & 2~min & 生成 summary、报告与可视化 \\
6 & 结果质检与必要复跑 & 5~min & 尺度量级、轴长、异常比例和级配无明显处理失败 \\
7 & 导出与记录 & 2~min & 保存场景图、参数、CSV、JSON、关键图与最终报告 \\
-- & 故障/重拍缓冲 & 6~min & 用于一次重拍、路径错误或模型重新加载 \\
\midrule
合计 & & 30~min & 完成从设备就位到可提交结果的闭环 \\
\bottomrule
\end{tabular}
\end{table}

现场操作应优先遵循“前端问题前端解决”的原则：曝光或模糊问题通过重拍解决，尺度问题通过重新标定解决，大面积遮挡通过重新布置允许调整的样品或更换视角解决；只有在输入质量合格后才运行固定模型。这样可以避免把有限时间浪费在对明显错误输入进行下游参数微调。

\subsection{自动化输出与结果追溯}
\label{sec:deployment-output}

每个场景建议建立独立输出目录，并至少保存四类文件。第一类是\textbf{原始输入证据}，包括原图和尺度参考信息；第二类是\textbf{实例证据}，包括 mask、composite/detection overview 以及逐颗粒几何 CSV；第三类是\textbf{物理统计结果}，包括筛分级配、$D_{10}/D_{50}/D_{90}$、形貌和异常的 JSON/CSV；第四类是\textbf{展示结果}，包括轴长叠加、级配对比图和最终文本/表格报告。四层文件共同构成从像素到提交数字的追溯链。

场景级 summary 中应记录至少以下元数据：scene id、resolution、模型版本、$\boldsymbol{\theta}$、$\gamma$、是否执行异常剔除、实例总数和正常实例数。正式加入机械筛分对照后，还应增加物理批次 id，使视觉数据与称量记录可以一一对应。若比赛需要现场展示，可由这些结构化文件自动生成一页摘要，而不是在现场手工复制数字到 PPT 或表格中。

自动输出还应保留\textbf{双口径}结果。数量型分布便于理解“可见颗粒个数”的视觉统计，质量代理型分布用于对齐机械筛分质量语义；二者同时保存可以帮助评委理解为什么当前方法不是简单的颗粒计数器。正式提交口径则应在赛前明确：最终与机械筛分比较时以经过独立校准并冻结的质量级配为主，数量分布作为辅助诊断。

\subsection{现场故障模式与恢复策略}
\label{sec:deployment-failure-recovery}

工程部署最常见的失败不一定来自模型本身，而可能来自文件、尺度或设备状态。表\ref{tab:field-recovery}列出需要在赛前演练的典型故障及恢复路径。恢复策略遵循一个原则：优先回到最靠近根因的阶段修复，而不是通过改变下游结果掩盖问题。

\begin{table}[htbp]
\centering
\caption{典型现场故障与恢复策略。}
\label{tab:field-recovery}
\small
\renewcommand{\arraystretch}{1.15}
\begin{tabular}{>{\raggedright\arraybackslash}p{3.2cm}>{\raggedright\arraybackslash}p{4.4cm}>{\raggedright\arraybackslash}p{5.5cm}}
\toprule
现象 & 可能根因 & 恢复策略 \\
\midrule
所有毫米粒径整体异常 & resolution 使用错误、参考物不共面、相机高度变化 & 停止下游解释，重新标定并重算几何/级配 \\
大片颗粒被合成一个实例 & 硬阴影、严重接触/遮挡、输入模糊 & 优先改善采集或允许时重新布置；再用固定配置重跑 \\
大量微小碎片实例 & 高光、纹理噪声、边缘破碎 & 检查原图与 mask；确认是分割残片还是实际细粒，不直接用阈值隐藏问题 \\
模型无法加载/GPU 不可用 & 权重路径错误、CUDA 环境异常 & 赛前缓存权重并做启动自检；必要时使用 CPU 作为低速后备但重新评估时间余量 \\
结果文件与场景不对应 & 手工路径/目录复用 & 每批唯一 scene id 和独立目录；自动写入参数与输入文件名 \\
$D_p$ 与肉眼预期差异很大 & 可能是真实级配差异，也可能是 merge、尺度或异常粗粒 & 按“尺度→实例→异常→统计”顺序排查，不以主观预期直接调 $\theta/\gamma$ \\
\bottomrule
\end{tabular}
\end{table}

经过上述设计，现场系统不再只是“运行两个 Python 命令”，而是形成了包含物理标定、输入质检、固定参数推理、中间结果复核、自动导出和故障恢复的工程闭环。下一章将在这一可操作流程基础上进一步总结方法目前仍然无法消除的观测边界和已知失败模式，并给出与这些根因直接对应的改进路线。
